% sections/01-introduction.tex

\section{Introduction}
\label{sec:intro}

% ── Motivating scenario ───────────────────────────────────────────────────────

\placeholder{Open with the motivating scenario: a surveillance system detects a
cluster in a region containing 3 reporting nodes. Publishing ``cluster detected in
region X'' effectively identifies which 3 nodes are affected. A data contributor who
knows their site is in region X and that the system detected a cluster there can infer
their own site is implicated. This is a disclosure risk the system creates at the
point of publication, regardless of any consent collected at the point of data
collection.}

Spatial disease surveillance systems routinely publish cluster detection alerts:
which geographic region exhibits a statistically significant excess of disease,
and how large that excess is.
These alerts are generated by the Kulldorff spatial scan statistic~\cite{kulldorff1997},
the standard method for this purpose, deployed through the SaTScan software by the
CDC, state health departments, and international agencies~\cite{desjardins2020,kulldorff2005}.

The key quantity in the scan statistic is the log-likelihood ratio $\LLR(w)$ over
a candidate window $w$, which aggregates case counts from $\kw$ reporting nodes
inside $w$.
When $\kw$ is small, one node's contribution dominates the aggregate, and publishing
the detected window's identity effectively discloses that node's data.
This is a privacy risk that arises \emph{at the point of publication}, not at the
point of data collection:
consent forms cannot cover downstream disclosure of which geographic cluster was
detected~\cite{persad2025,smith2025}.

% ── Why the problem is hard ───────────────────────────────────────────────────

Formalizing this risk requires analyzing a property the Kulldorff scan statistic
was never designed for.
The LLR formula~\cite{kulldorff1997} operates on aggregate counts and does not
expose $\kw$ as a parameter; the number of independent actors contributing to the
aggregate is irrelevant to detection but critical to privacy.
When one node changes its reported count by $\pm 1$, three quantities change
simultaneously: the observed count $\cobs$ inside the window, the national total
$\Ctot$, and the expected count $\muexp = \Ctot \cdot \nw / \Ntot$ for every
window globally.
No closed-form sensitivity analysis exists for this compound change.

% ── Gap statement ─────────────────────────────────────────────────────────────

Prior work has privatized chi-squared tests~\cite{gaboardi2016}, Gaussian likelihood
ratio tests~\cite{degue2018}, and F-statistics~\cite{alabi2022}, but the Kulldorff
spatial scan statistic has not been addressed.
Spatial privacy mechanisms~\cite{andres2013,machanavajjhala2008} and epidemiological DP
applications~\cite{fan2013} exist, but none analyzes the sensitivity of the LLR
or provides a private cluster selection mechanism for the scan statistic.
Heuristic spatial masking approaches~\cite{olson2006,cassa2006} quantify the
privacy-utility tradeoff empirically but provide no formal guarantees.

% ── Contributions ─────────────────────────────────────────────────────────────

\paragraph{Contributions.} This work makes three formal contributions:

\begin{enumerate}
    \item \textbf{Theorem~\ref{thm:sensitivity} (Sensitivity Bound):}
          A closed-form upper bound $\sens \leq f(\kw)$ on the global sensitivity
          of the Kulldorff LLR, where $f$ is monotonically decreasing in $\kw$
          and vanishes as $\kw \to \infty$.
          This is the first sensitivity analysis for the Kulldorff scan statistic.

    \item \textbf{Theorem~\ref{thm:mechanism} (Private Window Selection):}
          An instantiation of the exponential mechanism~\cite{mcsherry2007} with
          utility function $\util = \LLR(w)/\sens$, whose global sensitivity is
          $1$ by Theorem~\ref{thm:sensitivity}.
          The mechanism selects the privacy-adjusted optimal window (the candidate
          with the highest evidence-to-sensitivity ratio) with a formal probability
          guarantee in terms of the normalized utility gap $\gapu$.

    \item \textbf{Theorem~\ref{thm:composition} (Prospective Composition):}
          An application of zCDP composition~\cite{bun2016} to the prospective
          (weekly, repeated-release) surveillance setting, yielding a closed-form
          expression for the maximum operational duration $\Tmax$ within a fixed
          privacy budget.
\end{enumerate}

% ── Paper organization ────────────────────────────────────────────────────────

\paragraph{Organization.}
Section~\ref{sec:background} reviews the Kulldorff scan statistic and differential
privacy foundations.
Section~\ref{sec:method} presents the three theorems.
Section~\ref{sec:experiments} evaluates the mechanism on synthetic data.
Section~\ref{sec:conclusion} concludes.
